% Framework figure: (a) time grids, (b) shared encoder and layer heads,
% (c) stop-gradient compound loss.  Simplified from the proposal figures.
\begin{tikzpicture}[
  font=\footnotesize,
  >={Stealth[length=2mm]},
  plab/.style={font=\small\bfseries, text=cInk, anchor=west, inner sep=0pt},
  rowlab/.style={font=\scriptsize, text=cInk, anchor=east, inner sep=2pt},
  note/.style={font=\scriptsize, text=cInk!65},
  ebox/.style={draw=cInk!60, fill=cStop!10, rounded corners=2pt, align=center,
               inner sep=3pt, minimum height=6.5mm, text width=2.7cm, font=\scriptsize},
  ar/.style={->, draw=cInk!65, line width=0.6pt},
  hd/.style={->, draw=cStop!80!black, line width=0.8pt},
  lay/.style={draw=cInk!80, line width=1pt},
  dotn/.style={circle, fill=cInk, inner sep=1.3pt},
  rbox/.style={draw=cInk!55, fill=white, rounded corners=2pt, align=center,
               inner sep=3pt, font=\scriptsize},
  lbox/.style={draw=cInk!70, fill=cStop!10, rounded corners=2pt, align=center,
               inner sep=4pt, font=\scriptsize},
  grad/.style={->, draw=cJoint!85!black, line width=0.8pt, dash pattern=on 3pt off 1.5pt},
  gstop/.style={draw=cJoint!85!black, line width=0.8pt, dash pattern=on 3pt off 1.5pt,
                -{Bar[width=3mm,line width=1.4pt]}},
]

%======================================================================
% (a) Time grids
%======================================================================
\begin{scope}[shift={(0,0)}, x=0.31cm]
  \node[plab] at (-8.3,2.75) {(a) Time grids};
  \fill[cStop!12] (4,-0.12) rectangle (8,1.05);
  % fine grid
  \draw[->, cInk!70] (-0.5,0) -- (33,0);
  \foreach \n in {0,...,32}{\draw[cInk!45] (\n,-0.07) -- (\n,0.07);}
  \node[rowlab] at (-0.6,0) {fine grid $t_n$};
  % learning grid
  \draw[->, cInk!70] (-0.5,0.9) -- (33,0.9);
  \foreach \n in {0,4,...,32}{\draw[thick, cStop!80!black] (\n,0.77) -- (\n,1.03);}
  \foreach \n/\l in {0/0,4/1,8/2,12/3,16/4,20/5,24/6,28/7,32/8}{
    \node[font=\tiny, anchor=north, inner sep=1.5pt] at (\n,0.76) {$u_{\l}$};}
  \node[rowlab] at (-0.6,0.9) {learning grid $u_i$};
  % exercise dates
  \draw[->, cInk!70] (-0.5,1.8) -- (33,1.8);
  \foreach \n in {0,12,20,32}{\fill[cInk] (\n,1.8) circle (1.8pt);}
  \foreach \n/\l in {0/{T_0=0},12/{T_1},20/{T_2},32/{T_3=T}}{
    \node[font=\tiny, anchor=south, inner sep=2pt] at (\n,1.84) {$\l$};}
  \node[rowlab] at (-0.6,1.8) {exercise dates $T_j$};
  \foreach \n in {0,4,...,32}{\draw[densely dashed, cInk!20] (\n,0.08) -- (\n,0.76);}
  \foreach \n in {0,12,20,32}{\draw[densely dashed, cInk!20] (\n,1.04) -- (\n,1.76);}
  % one learning interval
  \draw[->, cStop!80!black, thick] (6,-0.45) -- (6,0.7);
  \node[note, anchor=west, text=cStop!70!black] at (6.3,-0.38)
    {$\ell_1=\LogSig^M(X)_{u_1,u_2}$, \quad $\Delta W_1=W_{u_2}-W_{u_1}$};
\end{scope}

%======================================================================
% (b) Shared encoder and layer heads
%======================================================================
\begin{scope}[shift={(0,-1.6)}]
  \node[plab] at (-2.57,0.75) {(b) Shared encoder and layer heads};
  \node[ebox] (E1) at (-0.6,-0.55) {fine-grid path $\Ipl S^{\mathrm f}$};
  \node[ebox] (E2) at (-0.6,-1.45) {local log-signatures $\ell_i$};
  \node[ebox] (E3) at (-0.6,-2.35) {shared NRDE \eqref{eq:logode}};
  \node[ebox] (E4) at (-0.6,-3.25) {hidden states $H_{u_i}$};
  \draw[ar] (E1) -- (E2); \draw[ar] (E2) -- (E3); \draw[ar] (E3) -- (E4);
  % heads bus
  \draw[cStop!80!black, line width=0.8pt] (1.55,-0.85) -- (1.55,-3.25);
  \draw[hd] (E4.east) -- (1.55,-3.25);
  \node[note, anchor=north, align=center] at (1.55,-3.4) {head input\\ $\xi_i=(H_{u_i},X_{u_i})$};
  % layers
  \foreach \y/\lab/\hl/\ta/\tb/\val in {%
     -0.85/{layer 1}/{y_0,\ \psi^Z_1}/{T_0}/{T_1}/{\hY_{1,T_1}},
     -1.85/{layer 2}/{\psi^Y_2,\ \psi^Z_2}/{T_1}/{T_2}/{\hY_{2,T_2}},
     -3.25/{layer $K$}/{\psi^Y_K,\ \psi^Z_K}/{T_{K-1}}/{T_K}/{\hY_{K,T}}}{
    \draw[hd] (1.55,\y) -- (3.6,\y);
    \node[font=\scriptsize, anchor=south, inner sep=1.5pt] at (2.55,\y) {$\hl$};
    \draw[lay] (3.65,\y) -- (9.2,\y);
    \foreach \x in {3.65,4.45,...,9.2}{\draw[cInk!50] (\x,\y-0.07) -- (\x,\y+0.07);}
    \node[dotn] at (3.65,\y) {}; \node[dotn] at (9.2,\y) {};
    \node[note, anchor=south west, inner sep=2pt] at (3.7,\y+0.03) {\lab};
    \node[font=\tiny, anchor=north, inner sep=2pt] at (3.65,\y-0.08) {$\ta$};
    \node[font=\tiny, anchor=north, inner sep=2pt] at (9.2,\y-0.08) {$\tb$};
    \node[font=\scriptsize, anchor=west, inner sep=3pt] at (9.3,\y) {$\val$};}
  \node at (6.4,-2.5) {$\vdots$};
  \draw[->, cInk!65, line width=0.8pt] (3.65,-0.2) -- (9.2,-0.2);
  \node[font=\scriptsize, anchor=south, inner sep=1.5pt] at (6.42,-0.2)
    {forward recursion \eqref{eq:scheme}: $\hY_{j,u_{i+1}}=\hY_{j,u_i}+\hZ_{u_i}^\top\Delta W_i$};
\end{scope}

%======================================================================
% (c) Stop-gradient compound loss
%======================================================================
\begin{scope}[shift={(0,-6.5)}]
  \node[plab] at (-2.57,0.25) {(c) Stop-gradient compound loss};
  % staircase of layers
  \draw[lay, ->] (-1.2,-0.3) -- (2.6,-0.3);
  \draw[lay, ->] (2.6,-1.0) -- (6.4,-1.0);
  \draw[lay, ->] (6.4,-1.7) -- (10.2,-1.7);
  \foreach \p in {(-1.2,-0.3),(2.6,-0.3),(2.6,-1.0),(6.4,-1.0),(6.4,-1.7),(10.2,-1.7)}{\node[dotn] at \p {};}
  \node[font=\scriptsize, anchor=east] at (-1.25,-0.3) {$y_0$};
  \node[font=\scriptsize, anchor=south east, inner sep=1pt] at (2.55,-0.25) {$\hY_{1,T_1}$};
  \node[font=\scriptsize, anchor=south east, inner sep=1pt] at (6.35,-0.95) {$\hY_{2,T_2}$};
  \node[font=\scriptsize, anchor=south, inner sep=2pt] at (10.2,-1.65) {$\hY_{K,T}$};
  \node[font=\scriptsize, anchor=north east, inner sep=1pt] at (2.55,-1.05) {$\hY_{2,T_1}$};
  \node[font=\scriptsize, anchor=north east, inner sep=1pt] at (6.35,-1.75) {$\hY_{3,T_2}$};
  % residual boxes
  \node[rbox] (e1) at (2.6,-2.75) {$e^{\sg}_1=\hY_{1,T_1}-\max\{\Phi^{\mathrm f}_1,\sg[\hY_{2,T_1}]\}$};
  \node[rbox] (e2) at (7.4,-2.75) {$e^{\sg}_2=\hY_{2,T_2}-\max\{\Phi^{\mathrm f}_2,\sg[\hY_{3,T_2}]\}$};
  \node[rbox] (eK) at (11.6,-2.75) {$e_K=\hY_{K,T}-\Phi^{\mathrm f}_K$};
  % layer 1 end value: forward and gradient
  \draw[ar] (2.5,-0.36) to[bend right=40] ($(e1.north)+(-1.6,0)$);
  \draw[grad] ($(e1.north)+(-1.9,0)$) to[bend left=40] (2.42,-0.3);
  \node[font=\tiny, text=cJoint!85!black, anchor=east, align=right] at (0.4,-1.2) {gradient\\ updates layer 1};
  % continuation value: forward, gradient blocked
  \draw[ar] (2.6,-1.06) -- ($(e1.north)+(0,0)$);
  \draw[gstop] ($(e1.north)+(0.35,0)$) -- ++(0,0.75);
  \node[font=\tiny, text=cJoint!85!black, anchor=west, align=left] at ($(e1.north)+(0.45,0.45)$) {gradient\\ blocked};
  \draw[ar] (6.4,-1.06) -- (6.4,-2.45);
  \draw[ar] (10.2,-1.76) -- ($(eK.north)+(-1.4,0)$);
  % loss
  \node[lbox] (L) at (5.0,-4.0) {$\cL_\omega(\theta)=\sum_{j=1}^K\omega_j\E\abs{e_j}^2$; \ one Adam step updates $y_0$, $\psi^Y_j$, $\psi^Z_j$ and the encoder};
  \draw[ar] (e1.south) -- (e1.south |- L.north);
  \draw[ar] (e2.south) -- (e2.south |- L.north);
  \draw[ar] (eK.south) |- (L.east);
\end{scope}
\end{tikzpicture}
